\documentclass[11pt, a4paper]{article}

\usepackage[a4paper, top=2.5cm, bottom=2.5cm, left=2cm, right=2cm]{geometry}
\usepackage{amsmath}
\usepackage{amssymb}
\usepackage[colorlinks=true, linkcolor=blue, urlcolor=blue, citecolor=blue]{hyperref}

\title{Theory of Everything - Volume 25 \\ \Large Black Hole Information Paradox}
\author{Omega Protocol}
\date{\today}

\begin{document}

\maketitle

\section*{Layman's Summary}
Stephen Hawking proved that black holes slowly evaporate and disappear. But if everything that fell into them is destroyed, this breaks the most sacred rule of physics: information cannot be destroyed. Volume 25 definitively solves the Black Hole Information Paradox, proving that black holes are the ultimate hard drives, not cosmic incinerators.

\section{Information Preservation (Axiom 28)}
We establish the strict preservation of quantum information as a fundamental axiom. Information that falls past the event horizon of a black hole is never deleted. Instead, the incredible gravitational forces instantly scramble the information and encode it onto the 2D surface of the horizon itself, like a holographic hard drive.

\section{The Page Curve (Theorem)}
If information isn't destroyed, how does it get out when the black hole evaporates? We prove that the escaping Hawking radiation eventually starts carrying the information out with it. The mathematical graph tracking this (the Page Curve) shows that once the black hole is half-evaporated, the outgoing radiation perfectly reconstructs the "destroyed" data.

\section{Firewall Resolution (Theorem)}
Physicists worried that preserving information would require a blazing wall of high-energy radiation (a "firewall") at the event horizon, burning anything that fell in. We resolve this paradox by proving that the interior of the black hole is simply a highly complex mathematical reconstruction of the exterior radiation. No firewall is needed; the math perfectly balances.

\end{document}
